%!TeX root=../emmatop.tex
\chapter[Chapter \thechapter]{}
\lettrine[lines=4,lraise=0.3]{I}{n} this state of schemes, and hopes, and connivance, June opened upon Hartfield. To Highbury in general it brought no material change. The Eltons were still talking of a visit from the Sucklings, and of the use to be made of their barouche-landau; and Jane Fairfax was still at her grandmother's; and as the return of the Campbells from Ireland was again delayed, and August, instead of Midsummer, fixed for it, she was likely to remain there full two months longer, provided at least she were able to defeat Mrs Elton's activity in her service, and save herself from being hurried into a delightful situation against her will.

Mr Knightley, who, for some reason best known to himself, had certainly taken an early dislike to Frank Churchill, was only growing to dislike him more. He began to suspect him of some double dealing in his pursuit of Emma. That Emma was his object appeared indisputable. Every thing declared it; his own attentions, his father's hints, his mother-in-law's guarded silence; it was all in unison; words, conduct, discretion, and indiscretion, told the same story. But while so many were devoting him to Emma, and Emma herself making him over to Harriet, Mr Knightley began to suspect him of some inclination to trifle with Jane Fairfax. He could not understand it; but there were symptoms of intelligence between them—he thought so at least—symptoms of admiration on his side, which, having once observed, he could not persuade himself to think entirely void of meaning, however he might wish to escape any of Emma's errors of imagination. She was not present when the suspicion first arose. He was dining with the Randalls family, and Jane, at the Eltons'; and he had seen a look, more than a single look, at Miss Fairfax, which, from the admirer of Miss Woodhouse, seemed somewhat out of place. When he was again in their company, he could not help remembering what he had seen; nor could he avoid observations which, unless it were like Cowper and his fire at twilight,

»Myself creating what I saw,«

brought him yet stronger suspicion of there being a something of private liking, of private understanding even, between Frank Churchill and Jane.

He had walked up one day after dinner, as he very often did, to spend his evening at Hartfield. Emma and Harriet were going to walk; he joined them; and, on returning, they fell in with a larger party, who, like themselves, judged it wisest to take their exercise early, as the weather threatened rain; Mr and Mrs Weston and their son, Miss Bates and her niece, who had accidentally met. They all united; and, on reaching Hartfield gates, Emma, who knew it was exactly the sort of visiting that would be welcome to her father, pressed them all to go in and drink tea with him. The Randalls party agreed to it immediately; and after a pretty long speech from Miss Bates, which few persons listened to, she also found it possible to accept dear Miss Woodhouse's most obliging invitation.

\begin{figure}[tbph]
\centering
\includegraphics[width=\linewidth]{41horseback}
\caption{Mr Perry passed by on horseback}
\end{figure}

As they were turning into the grounds, Mr Perry passed by on horseback. The gentlemen spoke of his horse.

»By the bye,« said Frank Churchill to Mrs Weston presently, »what became of Mr Perry's plan of setting up his carriage?«

Mrs Weston looked surprized, and said, »I did not know that he ever had any such plan.«

»Nay, I had it from you. You wrote me word of it three months ago.«

»Me! impossible!«

»Indeed you did. I remember it perfectly. You mentioned it as what was certainly to be very soon. Mrs Perry had told somebody, and was extremely happy about it. It was owing to her persuasion, as she thought his being out in bad weather did him a great deal of harm. You must remember it now?«

»Upon my word I never heard of it till this moment.«

»Never! really, never!—Bless me! how could it be?—Then I must have dreamt it—but I was completely persuaded—Miss Smith, you walk as if you were tired. You will not be sorry to find yourself at home.«

»What is this?—What is this?« cried Mr Weston, »about Perry and a carriage? Is Perry going to set up his carriage, Frank? I am glad he can afford it. You had it from himself, had you?«

»No, sir,« replied his son, laughing, »I seem to have had it from nobody.—Very odd!—I really was persuaded of Mrs Weston's having mentioned it in one of her letters to Enscombe, many weeks ago, with all these particulars—but as she declares she never heard a syllable of it before, of course it must have been a dream. I am a great dreamer. I dream of every body at Highbury when I am away—and when I have gone through my particular friends, then I begin dreaming of Mr and Mrs Perry.«

»It is odd though,« observed his father, »that you should have had such a regular connected dream about people whom it was not very likely you should be thinking of at Enscombe. Perry's setting up his carriage! and his wife's persuading him to it, out of care for his health—just what will happen, I have no doubt, some time or other; only a little premature. What an air of probability sometimes runs through a dream! And at others, what a heap of absurdities it is! Well, Frank, your dream certainly shews that Highbury is in your thoughts when you are absent. Emma, you are a great dreamer, I think?«

Emma was out of hearing. She had hurried on before her guests to prepare her father for their appearance, and was beyond the reach of Mr Weston's hint.

»Why, to own the truth,« cried Miss Bates, who had been trying in vain to be heard the last two minutes, »if I must speak on this subject, there is no denying that Mr Frank Churchill might have—I do not mean to say that he did not dream it—I am sure I have sometimes the oddest dreams in the world—but if I am questioned about it, I must acknowledge that there was such an idea last spring; for Mrs Perry herself mentioned it to my mother, and the Coles knew of it as well as ourselves—but it was quite a secret, known to nobody else, and only thought of about three days. Mrs Perry was very anxious that he should have a carriage, and came to my mother in great spirits one morning because she thought she had prevailed. Jane, don't you remember grandmama's telling us of it when we got home? I forget where we had been walking to—very likely to Randalls; yes, I think it was to Randalls. Mrs Perry was always particularly fond of my mother—indeed I do not know who is not—and she had mentioned it to her in confidence; she had no objection to her telling us, of course, but it was not to go beyond: and, from that day to this, I never mentioned it to a soul that I know of. At the same time, I will not positively answer for my having never dropt a hint, because I know I do sometimes pop out a thing before I am aware. I am a talker, you know; I am rather a talker; and now and then I have let a thing escape me which I should not. I am not like Jane; I wish I were. I will answer for it she never betrayed the least thing in the world. Where is she?—Oh! just behind. Perfectly remember Mrs Perry's coming.—Extraordinary dream, indeed!«

They were entering the hall. Mr Knightley's eyes had preceded Miss Bates's in a glance at Jane. From Frank Churchill's face, where he thought he saw confusion suppressed or laughed away, he had involuntarily turned to hers; but she was indeed behind, and too busy with her shawl. Mr Weston had walked in. The two other gentlemen waited at the door to let her pass. Mr Knightley suspected in Frank Churchill the determination of catching her eye—he seemed watching her intently—in vain, however, if it were so—Jane passed between them into the hall, and looked at neither.

There was no time for farther remark or explanation. The dream must be borne with, and Mr Knightley must take his seat with the rest round the large modern circular table which Emma had introduced at Hartfield, and which none but Emma could have had power to place there and persuade her father to use, instead of the small-sized Pembroke, on which two of his daily meals had, for forty years been crowded. Tea passed pleasantly, and nobody seemed in a hurry to move.

»Miss Woodhouse,« said Frank Churchill, after examining a table behind him, which he could reach as he sat, »have your nephews taken away their alphabets—their box of letters? It used to stand here. Where is it? This is a sort of dull-looking evening, that ought to be treated rather as winter than summer. We had great amusement with those letters one morning. I want to puzzle you again.«

Emma was pleased with the thought; and producing the box, the table was quickly scattered over with alphabets, which no one seemed so much disposed to employ as their two selves. They were rapidly forming words for each other, or for any body else who would be puzzled. The quietness of the game made it particularly eligible for Mr Woodhouse, who had often been distressed by the more animated sort, which Mr Weston had occasionally introduced, and who now sat happily occupied in lamenting, with tender melancholy, over the departure of the »poor little boys,« or in fondly pointing out, as he took up any stray letter near him, how beautifully Emma had written it.

Frank Churchill placed a word before Miss Fairfax. She gave a slight glance round the table, and applied herself to it. Frank was next to Emma, Jane opposite to them—and Mr Knightley so placed as to see them all; and it was his object to see as much as he could, with as little apparent observation. The word was discovered, and with a faint smile pushed away. If meant to be immediately mixed with the others, and buried from sight, she should have looked on the table instead of looking just across, for it was not mixed; and Harriet, eager after every fresh word, and finding out none, directly took it up, and fell to work. She was sitting by Mr Knightley, and turned to him for help. The word was blunder; and as Harriet exultingly proclaimed it, there was a blush on Jane's cheek which gave it a meaning not otherwise ostensible. Mr Knightley connected it with the dream; but how it could all be, was beyond his comprehension. How the delicacy, the discretion of his favourite could have been so lain asleep! He feared there must be some decided involvement. Disingenuousness and double dealing seemed to meet him at every turn. These letters were but the vehicle for gallantry and trick. It was a child's play, chosen to conceal a deeper game on Frank Churchill's part.

With great indignation did he continue to observe him; with great alarm and distrust, to observe also his two blinded companions. He saw a short word prepared for Emma, and given to her with a look sly and demure. He saw that Emma had soon made it out, and found it highly entertaining, though it was something which she judged it proper to appear to censure; for she said, »Nonsense! for shame!« He heard Frank Churchill next say, with a glance towards Jane, »I will give it to her—shall I?«—and as clearly heard Emma opposing it with eager laughing warmth. »No, no, you must not; you shall not, indeed.«

It was done however. This gallant young man, who seemed to love without feeling, and to recommend himself without complaisance, directly handed over the word to Miss Fairfax, and with a particular degree of sedate civility entreated her to study it. Mr Knightley's excessive curiosity to know what this word might be, made him seize every possible moment for darting his eye towards it, and it was not long before he saw it to be Dixon. Jane Fairfax's perception seemed to accompany his; her comprehension was certainly more equal to the covert meaning, the superior intelligence, of those five letters so arranged. She was evidently displeased; looked up, and seeing herself watched, blushed more deeply than he had ever perceived her, and saying only, »I did not know that proper names were allowed,« pushed away the letters with even an angry spirit, and looked resolved to be engaged by no other word that could be offered. Her face was averted from those who had made the attack, and turned towards her aunt.

»Aye, very true, my dear,« cried the latter, though Jane had not spoken a word—»I was just going to say the same thing. It is time for us to be going indeed. The evening is closing in, and grandmama will be looking for us. My dear sir, you are too obliging. We really must wish you good night.«

Jane's alertness in moving, proved her as ready as her aunt had preconceived. She was immediately up, and wanting to quit the table; but so many were also moving, that she could not get away; and Mr Knightley thought he saw another collection of letters anxiously pushed towards her, and resolutely swept away by her unexamined. She was afterwards looking for her shawl—Frank Churchill was looking also—it was growing dusk, and the room was in confusion; and how they parted, Mr Knightley could not tell.

He remained at Hartfield after all the rest, his thoughts full of what he had seen; so full, that when the candles came to assist his observations, he must—yes, he certainly must, as a friend—an anxious friend—give Emma some hint, ask her some question. He could not see her in a situation of such danger, without trying to preserve her. It was his duty.

»Pray, Emma,« said he, »may I ask in what lay the great amusement, the poignant sting of the last word given to you and Miss Fairfax? I saw the word, and am curious to know how it could be so very entertaining to the one, and so very distressing to the other.«

Emma was extremely confused. She could not endure to give him the true explanation; for though her suspicions were by no means removed, she was really ashamed of having ever imparted them.

»Oh!« she cried in evident embarrassment, »it all meant nothing; a mere joke among ourselves.«

»The joke,« he replied gravely, »seemed confined to you and Mr Churchill.«

He had hoped she would speak again, but she did not. She would rather busy herself about any thing than speak. He sat a little while in doubt. A variety of evils crossed his mind. Interference—fruitless interference. Emma's confusion, and the acknowledged intimacy, seemed to declare her affection engaged. Yet he would speak. He owed it to her, to risk any thing that might be involved in an unwelcome interference, rather than her welfare; to encounter any thing, rather than the remembrance of neglect in such a cause.

»My dear Emma,« said he at last, with earnest kindness, »do you think you perfectly understand the degree of acquaintance between the gentleman and lady we have been speaking of?«

»Between Mr Frank Churchill and Miss Fairfax? Oh! yes, perfectly.—Why do you make a doubt of it?«

»Have you never at any time had reason to think that he admired her, or that she admired him?«

»Never, never!« she cried with a most open eagerness—»Never, for the twentieth part of a moment, did such an idea occur to me. And how could it possibly come into your head?«

»I have lately imagined that I saw symptoms of attachment between them—certain expressive looks, which I did not believe meant to be public.«

»Oh! you amuse me excessively. I am delighted to find that you can vouchsafe to let your imagination wander—but it will not do—very sorry to check you in your first essay—but indeed it will not do. There is no admiration between them, I do assure you; and the appearances which have caught you, have arisen from some peculiar circumstances—feelings rather of a totally different nature—it is impossible exactly to explain:—there is a good deal of nonsense in it—but the part which is capable of being communicated, which is sense, is, that they are as far from any attachment or admiration for one another, as any two beings in the world can be. That is, I presume it to be so on her side, and I can answer for its being so on his. I will answer for the gentleman's indifference.«

She spoke with a confidence which staggered, with a satisfaction which silenced, Mr Knightley. She was in gay spirits, and would have prolonged the conversation, wanting to hear the particulars of his suspicions, every look described, and all the wheres and hows of a circumstance which highly entertained her: but his gaiety did not meet hers. He found he could not be useful, and his feelings were too much irritated for talking. That he might not be irritated into an absolute fever, by the fire which Mr Woodhouse's tender habits required almost every evening throughout the year, he soon afterwards took a hasty leave, and walked home to the coolness and solitude of Donwell Abbey.